Hij keek in haar ogen.
Door glans aangezogen,
zonder mededogen,

oogstte hij haar treurnis.
Ze zocht vergiffenis
want de roofwildernis

in zijn lubrieke blik
en zijn plagend getik
zette haar in de fik.

Ze voelde verbinding
via haar opwinding
want zijn wijsvinger ging

steeds harder over haar
jeukende kittelaar.
Ze voelde zich kwetsbaar.

Oogcontact al die tijd
praamde intimiteit
door stijve nabijheid.

Het communiceren
zonder te oreren
deed romantiseren.

Als Emmy wegdraaide
omdat lust oplaaide
als zijn hand haar naaide,

moest hij haar afzeiken.
“Blijf me steeds aankijken
tijdens mijn aanstrijken.

Je ogen vertellen
dat mijn vingers kwellen.
Laat je klit opzwellen

en toon me je genot.
Ik zie ook het verbod
dat je misschien van God

meent te krijgen voor dit
gepriegel aan je klit.
Ik geniet van de split

die ik tussen lichaam
en je geest onbuigzaam
zie ontstaan, heel langzaam.”

Lalde hij ladderzat.
Maar dan minder gevat
en in het Grunnegs Platt.

Emmy kreeg empathie,
maar geen compassie
voor deze perceptie.

Emmy ervoer handen,
die zowel aanrandden
als haar lieten branden.

Dat werd een ezelsbrug.
Het bracht haar brein terug
naar een man op zijn rug.

Toen waren ze van haar.
Toen had ze geen bezwaar.
Toen hoopte ze zowaar

aaiend over zijn borst
op een harde bloedworst
van de nog jonge vorst.

De jongeheer van stand
stuurde haar zachte hand
kalm naar zijn onderkant.

Ze lagen naast een boo
vooral gemaakt van stro.
Hij dacht: "Het mag wel zo!”

Emmy haalde van stal,
bij deze koeienstal,
wat ze veel eerder al

grondig had bestudeerd.
Ze had dit gestileerd
van de pesters geleerd.

Zijn verlangen naar meer
bleef bij gefantaseer.
Hij bleek een echte heer.

Ze waren niet gehuwd.
Dus werd de lust gestuwd.
Emmy was gewaarschuwd

niet te worden verleid
om haar maagdelijkheid
in de verkeringstijd

al te gaan verbeuren.
“Hoe hij ook zal zeuren,
laat dat niet gebeuren.”

Het bleef bij zinspeling.
Doch door die voorstelling
steeg al de opwinding.

Tijdens Emmy’s streling,
die over zijn kruis ging,
kwam de herinnering

aan al die lulletjes
met al die krulletjes
van al die knulletjes.

Ze had goed opgelet
welke rukkende zet
hun handen had besmet.

“Ontspan en doe het niet”,
zei ze, zoals het lied
“Relax” dat ook gebiedt.

Hij had porno ontmoet
niet ver van zijn landgoed,
in plaats van Hollywood.

Hij ging niet naar Frankie,
laat staan een Zweeds meissie,
maar kreeg het bij Emmy.

Westervördication
werd het, in plaats van een
Californication.

Nu ze waren verloofd,
leek haar wel geoorloofd
dat haar hand zijn zaad rooft

en hem voldoening schenkt.
“Gedraag je je gekrenkt
omdat je nog steeds denkt

dat ik je slinks bedrieg?
Leef als een eendagsvlieg!”,
bleek onbedoeld gelieg.

De ringen wisselen
bleek geen bedisselen.
Emmy niet inwisselen

voor een mooie premie
gaf hem satisfactie.
Wat was dan met Emmy

dat zijn bevrediging
leidde tot tuchtiging
of zelfs vernietiging?

Wat kostte dit moment,
waarin hij werd verwend,
de jonkheer aan het end?

De jongeheer van stand
stond al volop in brand
toen Emmy’s kleine hand

in zijn kniebroek verdween.
Haar lover kwam meteen.
Ze keek vlug om haar heen.

Ver op de achtergrond
blafte een herdershond.
Hoewel ze niets verstond,

zag ze uit de bossen,
rijdend op hun rossen,
mannen woorden tossen.

Hun stemmen klonken kwaad.
Terwijl Emmy zijn zaad
in haar rechterhand laadt,

kwamen zij die strijden
hun richting oprijden.
Om de troep te mijden,

vluchtte Emmy meteen.
Haar vriend liet ze alleen.
Dat was geenszins gemeen,

maar uit medeleven.
Een vogelvrij leven
zou hem kommer geven.

Ze zou met hem paren,
zijn kinderen baren,
maar dit leed besparen.

De jongeheer in shock,
toen Emmy plots vertrok
nadat ze hem aftrok,

keek haar verbijsterd na.
Ze wuifde hem daarna
met zijn eigen sperma,

dat ze eerder opving
als een mooi hebbeding
wat naar hem terug ging.

Het landde op zijn hoofd,
verbijsterd en verdoofd.
Hij had haar echt geloofd.

“Met deze ontsnapping
creëer ik zijn redding,
gezien mijn ontvluchting”,

zei haar scherpe verstand.
Ze creëerde afstand.
Zwom naar de overkant

van een brede dijksloot,
die de zekerheid bood
om aan een snelle dood

even te ontkomen.
De prins uit haar dromen,
nog niet bijgekomen,

zonet gom gespoten,
had tot strijd besloten.
Emmy hoorde schoten.

Zo verloor ze haar lief,
haar eerste hartedief,
na slechts een kort gerief.

De liefdesmachine
bleek als Melusine.
Eerst oxytocine,
alsook endorfine,
met wat dopamine.
Daarna prolactine.

Tot hij ontdekte, pas
in het hoge veengras,
wie ze in wezen was.

Maar Emmy’s vijanden
bleken niet te stranden
toen ze ook belandden

nabij het veenkanaal.
Ze vonden een gemaal
dat de clan allemaal

een oversteek aanbood.
Emmy bleef dus in nood.
Ze kreeg hulp van een Jood,

die zijn keppel afdeed
en in zijn huifkar smeet.
Hij greep Emmy toen beet,

die eerst wilde vluchten,
maar niets had te duchten.
In plaats van bevruchten,

verborg hij haar dienstbaar
onder zijn handelswaar
met een stiltegebaar

en eentje van “sterkte.”
De strategie werkte.
Niemand van hen merkte

de verborgene op.
De kar kwam tot een stop.
Hij riskeerde de strop,

maar had Emmy gered.
Hij herpakte zijn pet
toen hij haar had ontzet.

Door haar vooroordelen
kon ze niet klaarspelen
haar dank mee te delen.

Voordat ze bang wegging,
griste ze wat kleding
uit zijn wagenlading.

Eenmaal in veiligheid
kwam er een scherpe spijt
over haar dommigheid.

Ze had haar verloofde,
die in haar geloofde
en zijn trouw beloofde,

in de steek gelaten.
Haar spijt mocht niet baten!
Wel het zichzelf haten?
